\documentclass[12pt]{article}
\usepackage[T1]{fontenc}
\usepackage[utf8]{inputenc}
\usepackage{lmodern}
\usepackage{textcomp}
\usepackage{enumitem}
\usepackage{geometry}
\geometry{margin=1in}
\begin{document}
\begin{center}
Answer Key - Religious Studies - K-12 Level Assessment 22 \
\small (Answers are ordered exactly as in the question paper.)
\end{center}

\section*{Multiple Choice}
\begin{enumerate}[label=\arabic*.]
\item \textbf{Answer:} B
\\ \textit{Options:} A) Es; B) Izanagi; C) Izanami; D) Kami
\\ \textit{Note:} Izanagi is the deity in Japanese mythology who, after returning from the underworld, cleanses himself and gives birth to various gods, including Amaterasu, the sun goddess, from his left eye.
\item \textbf{Answer:} B
\\ \textit{Options:} A) Epistemology; B) Morality; C) Religion; D) Aesthetics
\\ \textit{Note:} In his final work, Laws, Plato focused on morality as he sought to establish a framework for ethical governance and the moral education of citizens, moving away from his earlier cosmological inquiries.
\item \textbf{Answer:} A
\\ \textit{Options:} A) Kali yuga; B) Maha yuga; C) Krita yuga; D) Satya yuga
\\ \textit{Note:} According to Hindu tradition, we are currently in the Kali yuga, which is the last of the four stages of the world cycle, characterized by a decline in virtue and an increase in chaos and moral deterioration.
\item \textbf{Answer:} A
\\ \textit{Options:} A) Baptism and Eucharist; B) Baptism and Ordination; C) Eucharist and Marriage; D) Eucharist and Confession
\\ \textit{Note:} Baptism and Eucharist are considered the two sacraments universally recognized by all Christian denominations, symbolizing initiation into the faith and the remembrance of Christ's sacrifice, respectively.
\item \textbf{Answer:} C
\\ \textit{Options:} A) Guan-yin; B) Kannon; C) Ojizo-sama; D) Amitabha
\\ \textit{Note:} Ojizo-sama is revered in Japanese Buddhism as the bodhisattva who guides the souls of deceased children to the afterlife, providing them with protection and comfort during their journey.
\end{enumerate}

\section*{True / False}
\begin{enumerate}[label=\arabic*.]
\setcounter{enumi}{5}
\item \textbf{Answer:} False
\\ \textit{Options:} A) True; B) False
\\ \textit{Note:} Hanukkah is not referred to as the Festival of Booths; that title belongs to Sukkot, which celebrates the harvest and commemorates the Israelites' journey in the desert.
\item \textbf{Answer:} True
\\ \textit{Options:} A) True; B) False
\\ \textit{Note:} Humanistic Buddhism promotes the idea that laypeople can actively engage in Buddhist practices and social issues, emphasizing the importance of applying Buddhist principles to everyday life and community welfare.
\item \textbf{Answer:} False
\\ \textit{Options:} A) True; B) False
\\ \textit{Note:} Tibetan Buddhism was not the first form of Buddhism to gain a foothold in North America; that distinction belongs to Zen Buddhism, which began to attract attention in the mid-20 th century.
\item \textbf{Answer:} False
\\ \textit{Options:} A) True; B) False
\\ \textit{Note:} The Isma'ilis are a branch of Shia Islam, not Sunni, and therefore do not belong to the largest sect of Islam.
\item \textbf{Answer:} True
\\ \textit{Options:} A) True; B) False
\\ \textit{Note:} The Qur'an permits polygyny, specifically in Surah An-Nisa (4:3), which allows a man to marry up to four wives if he can treat them justly and provide for them financially, reflecting the socio-economic context of the time.
\end{enumerate}

\section*{Long Form Response}
\begin{enumerate}[label=\arabic*.]
\setcounter{enumi}{10}
\item \textbf{Answer:} Thich Nhat Hanh emphasizes the theme of self-sacrifice in his teachings by highlighting the importance of compassion and mindfulness in our relationships with others. He believes that true happiness comes from serving others and being present in their lives, which often requires letting go of our own desires and comforts. This self-sacrifice is significant within Buddhist philosophy as it aligns with the concepts of interdependence and the idea that our well-being is connected to the well-being of others. By practicing self-sacrifice, individuals cultivate a deeper sense of community and peace, reflecting the core Buddhist teachings of loving-kindness and altruism. Through his teachings, Thich Nhat Hanh inspires followers to engage in acts of selflessness, fostering a more harmonious and compassionate world.
\\ \textit{Note:} correct because it accurately connects Thich Nhat Hanh's teachings on self-sacrifice to the broader Buddhist principles of compassion, mindfulness, and interdependence, illustrating how individual happiness is intertwined with the well-being of others.
\item \textbf{Answer:} Antiochus IV Epiphanes significantly impacted the Maccabean Revolt by enforcing Hellenization and suppressing Jewish religious practices, which triggered a strong backlash among the Jewish people. Politically, his imposition of Greek culture and religion, including the desecration of the Temple in Jerusalem, was seen as an affront to Jewish identity and faith. This led to the rise of the Maccabees, a group of Jewish rebels who fought to reclaim their religious freedom and autonomy. The revolt ultimately resulted in the rededication of the Temple and the establishment of the Hasmonean dynasty, marking a crucial moment in Jewish history that reinforced the importance of religious identity and resistance against oppression.
\\ \textit{Note:} correct because it identifies Antiochus IV Epiphanes' enforcement of Hellenization and suppression of Jewish practices as key triggers for the Maccabean Revolt, highlighting the conflict's religious and political dimensions and its significance in affirming Jewish identity and autonomy.
\end{enumerate}

\vfill
\noindent\textit{End of Answer Key}
\end{document}